\section{Source Code}
\label{sec:source-code}

\subsection{DataExtraction.py}
\begin{lstlisting}[language=Python,label={lst:data-extraction}]
import pandas as pd

data_text = """
Angle: 135°; Distance: 218.58 cm
Angle: 134°; Distance: 274.31 cm
Angle: 133°; Distance: 0.00 cm
Angle: 132°; Distance: 0.00 cm
Angle: 131°; Distance: 188.98 cm
... (from Arduino IDE)
Angle: 49°; Distance: 115.88 cm
Angle: 48°; Distance: 197.79 cm
Angle: 47°; Distance: 275.62 cm
Angle: 46°; Distance: 0.00 cm
Angle: 45°; Distance: 0.00 cm
"""

data_lines = data_text.strip().split("\n")
extracted_data = []

for line in data_lines:
    parts = line.split(";")
    angle_part = parts[0].strip()
    distance_part = parts[1].strip()

    angle = int(angle_part.split(":")[1].strip().replace("°", ""))
    distance = float(distance_part.split(":")[1].strip().replace("cm", "").strip())

    extracted_data.append({"Angle": angle, "Distance / cm": distance})

df = pd.DataFrame(extracted_data)

df.to_excel("./DataProcessing/360cm/360cm_data_03.xlsx", index=False)
\end{lstlisting}

\subsection{DataAverage.py}
\begin{lstlisting}[language=Python,label={lst:data-average}]
import pandas as pd

file_paths = [
    "./DataProcessing/280cm/280cm_data_01.xlsx",
    "./DataProcessing/280cm/280cm_data_02.xlsx",
    "./DataProcessing/280cm/280cm_data_03.xlsx"
]

data_frames = [pd.read_excel(file) for file in file_paths]

combined_data = pd.concat(data_frames)
average_data = combined_data.groupby(combined_data.index).mean()

average_data = average_data.round(2)

output_file = "./DataAverage/280cm_average_data.xlsx"
average_data.to_excel(output_file, index=False)
\end{lstlisting}

\subsection{Black\&White.py}
\begin{lstlisting}[language=Python,label={lst:black-and-white}]
import pandas as pd
import numpy as np
import matplotlib.pyplot as plt

file_paths = [f"./DataRestoration/{i}cm_0&1_data.xlsx" for i in range(0, 401, 40)]

distance_steps = np.arange(0, 401, 40)
angle_steps = np.arange(45, 136)
data_matrix = np.zeros((len(distance_steps), len(angle_steps)))

for idx, file_path in enumerate(file_paths):
    df = pd.read_excel(file_path)
    for _, row in df.iterrows():
        angle = int(row["Angle"])
        distance = row["Distance / cm"]
        if 45 <= angle <= 135:
            angle_idx = angle - 45
            data_matrix[idx, angle_idx] = distance

binary_matrix = np.where(data_matrix > 0, 1, 0)

angles = np.deg2rad(angle_steps)
distances = distance_steps
theta, r = np.meshgrid(angles, distances)

fig, ax = plt.subplots(subplot_kw={'projection': 'polar'}, figsize=(8, 8))

cmap = plt.cm.binary

heatmap = ax.pcolormesh(theta, r, binary_matrix, cmap=cmap, shading='auto')

ax.set_title("Radar Graph")

ax.set_thetamin(45)
ax.set_thetamax(135)
ax.set_rticks(distance_steps)
ax.set_rlabel_position(90)

plt.savefig("./Photos/RadarGraph.png", dpi=1200)
plt.show()
\end{lstlisting}

\subsection{GraphGenerate.py}
\begin{lstlisting}[language=Python,label={lst:graph-generate}]
import numpy as np
import pandas as pd
import matplotlib.pyplot as plt

file_name = "./DataAverage/280cm_average_data.xlsx"

df = pd.read_excel(file_name)
angles = df["Angle"].values
distances = df["Distance / cm"].values

valid_indices = distances > 0
angles = angles[valid_indices]
distances = distances[valid_indices]

theta = np.deg2rad(angles)

fig, ax = plt.subplots(subplot_kw={'projection': 'polar'}, figsize=(8, 8))

cmap = plt.cm.YlOrRd
vmin = 0
vmax = max(distances)
norm = plt.Normalize(vmin=vmin, vmax=vmax)
colors = cmap(norm(distances))

for i in range(len(theta)):
    ax.bar(
        theta[i],
        distances[i],
        width=np.deg2rad(1),
        color=colors[i],
        edgecolor="none",
        alpha=1.0
    )

sm = plt.cm.ScalarMappable(cmap=cmap, norm=norm)
sm.set_array([])

ax.set_title("280cm", va="bottom", fontsize=14)
ax.set_thetamin(45)
ax.set_thetamax(135)
ax.set_theta_zero_location("E")

plt.savefig("./Photos/angle-distance/280cm.png", dpi=1200)
plt.show()
\end{lstlisting}

\subsection{HeatmapGenerate.py}
\begin{lstlisting}[language=Python,label={lst:heatmap-generate}]
import pandas as pd
import mlx.core as mx
import matplotlib.pyplot as plt

file_paths = [f"./DataAverage/{i}cm_average_data.xlsx" for i in range(0, 401, 40)]

distance_steps = mx.arange(0, 400, 40)
angle_steps = mx.arange(45, 136)
data_matrix = mx.zeros((len(distance_steps), len(angle_steps)))

for idx, file_path in enumerate(file_paths):
    df = pd.read_excel(file_path)
    for _, row in df.iterrows():
        angle = int(row["Angle"])
        distance = row["Distance / cm"]
        if 45 <= angle <= 135:
            angle_idx = angle - 45
            data_matrix[idx, angle_idx] = distance

angles = mx.degrees(angle_steps)
distances = distance_steps
theta, r = mx.meshgrid(angles, distances)

fig, ax = plt.subplots(subplot_kw={'projection': 'polar'}, figsize=(8, 8))
cmap = plt.cm.Oranges
heatmap = ax.pcolor(theta, r, data_matrix, cmap=cmap, shading='auto')

cbar = plt.colorbar(
    heatmap,
    ax=ax,
    orientation='vertical',
    pad=0.1,
    shrink=0.6
)
cbar.set_label("Distance / cm", fontsize=12)

ax.set_title("Heatmap of \"Angle - Distance\"", va="bottom", fontsize=14)
ax.set_thetamin(45)
ax.set_thetamax(135)
ax.set_rticks(distance_steps)
ax.set_rlabel_position(90)

plt.savefig("./Photos/heatmap.png", dpi=1200)
plt.show()
\end{lstlisting}

\subsection{distanceDetect.ino}
\begin{lstlisting}[language=c++,label={lst:distance-detect}]
#include <Servo.h>
Servo engine;

// steering gear
int enginePin = 3;

// ultrasonic sensor
int ultrasound_TrigPin = 10;
int ultrasound_EchoPin = 11;

int counter = 0;

void setup() {
  // put your setup code here, to run once:
  Serial.begin(9600);
  // Engine Driver
  engine.attach(enginePin);
  // Ultrasound Driver
  pinMode(ultrasound_TrigPin, OUTPUT);
  pinMode(ultrasound_EchoPin, INPUT);
}

void loop() {
  // put your main code here, to run repeatedly:
  if(counter < 6) {
    engineTurn();
  } else if(engine.attached()) {
    engine.detach();
    Serial.println("STOP!!!");
  }
}

int angle = 90;
int step = 1;
unsigned long previousMillis = 0;

void engineTurn() {
  unsigned long currentMillis = millis();

  if(currentMillis - previousMillis >= 20) {
    previousMillis = currentMillis;

    engine.write(angle);

    ultrasoundControl();

    angle += step;

    // Change the direction when arrive the boundary
    if(angle >= 135) {
      angle = 135;
      step = -step;
      counter++;
    } else if(angle <= 45) {
      angle = 45;
      step = -step;
      counter++;
    }
  }
}

float duration, distance;

void ultrasoundControl() {
  digitalWrite(ultrasound_TrigPin, LOW);
  delayMicroseconds(2);
  digitalWrite(ultrasound_TrigPin, HIGH);
  delayMicroseconds(10);
  digitalWrite(ultrasound_TrigPin, LOW);

  duration = pulseIn(ultrasound_EchoPin, HIGH, 25000);
  distance = (duration * 0.0343) / 2;

  Serial.print("Angle: ");
  Serial.print(angle);
  Serial.print("°; Distance: ");
  Serial.print(distance);
  Serial.println(" cm");
}
\end{lstlisting}
